\section{A Bessel limit for the smallest roots}
\label{sec:bessel-edge}

The fixed largest roots have size $n/j$, while the bulk law leaves the
smallest roots unresolved. The opposite endpoint has a different,
explicit scaling limit. It follows directly from the same Stirling
coefficient formula, and completes the leading description of the
three root regimes considered here.

For $n,k\geq1$, use the harmless scalar normalization
\begin{equation}\label{eq:bessel-normalization}
 P_{n,k}(X)=\frac{(-1)^{n+k+1}L_{n,k}(X)}
                   {k!\binom{n+k-1}{n}}
 =\sum_{j=1}^n(-1)^j\stirling{n}{n-j+1}
                     \frac{X^{k-1+j}}{(k-1+j)!}.
\end{equation}
In particular $P_{n,1}=p_n$. Let $z_{k,j}$ be the $j$th positive
zero of the Bessel function $J_k$; these zeros are simple
\cite{DLMFBessel}.

\begin{theorem}[An entire Bessel scaling limit]
\label{thm:bessel-edge}
For each fixed positive integer $k$, locally uniformly for complex $x$,
\begin{equation}\label{eq:bessel-limit}
 n^{2k}P_{n,k}(x/n^2)
       \longrightarrow -(2x)^{k/2}J_k(\sqrt{2x}).
\end{equation}
The right side is understood by its entire power series, so it has
no square-root branch ambiguity. If $x_{n,k,j}$ denotes the $j$th
smallest positive root of $L_{n,k}$, then for every fixed $j,k\geq1$,
\begin{equation}\label{eq:small-root-limit}
 n^2x_{n,k,j}\longrightarrow \frac{z_{k,j}^2}{2}.
\end{equation}
For the diagonal zero-velocity lattice parameters of
\cref{sec:velocity-zero-law}, this says
\begin{equation}\label{eq:small-spacing-limit}
 n^2\log A_{n,j}\longrightarrow \frac{z_{1,j}^2}{2},
 \qquad n^2(A_{n,j}-1)\longrightarrow\frac{z_{1,j}^2}{2}.
\end{equation}
\end{theorem}
\begin{proof}
Write $m=j-1$. The standard product for unsigned Stirling numbers gives
\[
 \stirling{n}{n-m}=e_m(1,2,\ldots,n-1),
\]
where $e_m$ is the elementary symmetric function of degree $m$.
For each fixed $m$,
\begin{equation}\label{eq:stirling-fixed-codimension}
 n^{-2m}\stirling{n}{n-m}
           \longrightarrow \frac1{2^m m!}.
\end{equation}
Indeed $(\sum_{r=1}^{n-1}r)^m$ is $m!e_m$ plus terms with repeated
indices. For $m\geq2$ the sum of the latter is bounded by
\[
 \binom m2\left(\sum_{r=1}^{n-1}r^2\right)
          \left(\sum_{r=1}^{n-1}r\right)^{m-2}
       =O_m(n^{2m-1}).
\]
The cases $m=0,1$ are direct, and
$n^{-2}\sum_{r=1}^{n-1}r\to1/2$.
For all $m,n$ the same positive expansion gives the useful bound
\begin{equation}\label{eq:stirling-domination}
 0\leq n^{-2m}\stirling{n}{n-m}
       \leq\frac1{2^m m!},
\end{equation}
with the coefficient interpreted as zero for $m\geq n$.

Substituting \eqref{eq:bessel-normalization} yields
\[
 n^{2k}P_{n,k}(x/n^2)
  =-\sum_{m=0}^{n-1}(-1)^m
       n^{-2m}\stirling{n}{n-m}\frac{x^{k+m}}{(k+m)!}.
\]
On every compact set, \eqref{eq:stirling-domination} gives an
absolutely summable common bound for its terms. Termwise passage to
the limit proves
\[
 -\sum_{m\geq0}\frac{(-1)^m x^{k+m}}{2^m m!(k+m)!}
       =-(2x)^{k/2}J_k(\sqrt{2x}),
\]
the defining Bessel series \cite{DLMFBessel}, and proves local uniformity.

The limiting function has a zero of multiplicity $k$ at zero and
simple positive zeros $z_{k,j}^2/2$. Apply Rouch\'e's theorem on
disjoint small circles around zero and the first finitely many positive
zeros. The polynomial on the left has exactly multiplicity $k$ at
zero, and its remaining zeros are real, simple, and positive by the
operator theorem in \cref{sec:cohen-edges}. Uniform nonvanishing on
the compact real intervals between those circles therefore identifies
the first positive polynomial roots with the respective limiting
zeros; no additional small root can enter at zero. This proves
\eqref{eq:small-root-limit}. For $k=1$, $x_{n,1,j}=\log A_{n,j}$.
Since this tends to zero, $e^x-1\sim x$ gives
\eqref{eq:small-spacing-limit}.
\end{proof}

For example, $z_{1,1}^2/2=7.340985321\ldots$. This endpoint constant
is not obtained by inserting a fixed $j$ into the bulk quantile formula:
the bulk theorem assumes $j/(n-1)$ stays away from both endpoints.
A uniform bridge between the Bessel regime, the angle-distributed
bulk, and the largest-root expansion remains a substantive further
problem.
